\section{The identity component of the Picard scheme}
\label{mk-chap:Picard_IdentityComponent}
\providecommand{\connComp}{\operatorname{Comp}}

\section{The identity component of a group scheme}
\label{mk-sec:identity_component_abstract}


\begin{definition}[Identity component of a group scheme]\label{mk-def:identity_component_group_scheme}
  \dcref{5.1,custom}

  \uses{mk-def:identity_component_carrier}
  \textit{Source: \cite{kleiman-picard}, ``The Picard scheme'', \S 5, Lem.~5.1.}
  Let \(k\) be a field and \(G\) a \(k\)-group scheme locally of finite type
  whose identity section \(e : \Spec k \to G\) factors through \(G\). The
  \emph{identity component} of \(G\), denoted \(G^0\), is the locally closed
  subscheme of \(G\) whose underlying topological space is the connected
  component of \(|G|\) containing the image of \(e\), equipped with the
  reduced-induced (a priori) subscheme structure inherited from \(G\).
  In characteristic zero, or more generally when \(G\) has a
  geometrically reduced open subscheme, the substructure agrees with
  the open subscheme structure of \ref{mk-thm:identity_component_open_subgroup}.
\end{definition}

\begin{theorem}
[Identity component is a clopen subscheme]
  \label{mk-thm:identity_component_open_subgroup}
  \dcref{5.1,custom}

  \textit{Source: \cite{kleiman-picard}, ``The Picard scheme'', \S 5, Lem.~5.1~(3).}
  Let \(k\) be a field and \(G\) a \(k\)-group scheme locally of finite
  type. Then the identity component \(G^0\)
  (\ref{mk-def:identity_component_group_scheme}) is an open and closed
  subscheme of \(G\): the inclusion \(G^0 \hookrightarrow G\) is both an
  open immersion and a closed immersion of schemes.
\end{theorem}
\begin{proof}
  \uses{mk-def:identity_component_group_scheme, mk-def:identity_component_carrier, mk-lem:identity_component_locally_connected}
  Following Kleiman's proof of Lem.~5.1~(3) (clopen-subscheme part).
  Since \(e\) is a \(k\)-rational point, the singleton \(\{e\}\) is a closed
  point of \(G\). The connected component \(G^0\) of \(|G|\) at \(e\) is the
  largest connected subspace containing \(e\); the closure of a connected
  subspace is connected, hence \(G^0\) is closed. On the other hand, a
  locally Noetherian topological space has open connected components
  (EGA I 6.1.9), so \(G^0 \subseteq |G|\) is also open. Equip \(G^0\) with
  the open-subscheme structure; the inclusion
  \(G^0 \hookrightarrow G\) is then both an open and a closed immersion.
\end{proof}

\begin{theorem}[Identity component inclusion is a group-scheme homomorphism]\label{mk-thm:identity_component_is_subgroup_homomorphism}
  \dcref{5.1,custom}

  \uses{mk-thm:identity_component_open_subgroup, mk-def:identity_component_group_scheme, mk-lem:identity_component_geometrically_connected}
  \textit{Source: \cite{kleiman-picard}, ``The Picard scheme'', \S 5,
  Lem.~5.1~(3); the verbatim source quote of (3) is at
  \ref{mk-thm:identity_component_open_subgroup}.}
  Let \(k\) be a field and \(G\) a \(k\)-group scheme locally of finite
  type. The clopen inclusion \(G^0 \hookrightarrow G\)
  (\ref{mk-thm:identity_component_open_subgroup}) is a homomorphism of
  \(k\)-group schemes: \(G^0\) inherits a \(k\)-group-scheme structure
  from \(G\), and the inclusion morphism in the over-category over
  \(\Spec k\) is compatible with the group laws on source and target.
\end{theorem}

\begin{proof}
  \uses{mk-thm:identity_component_open_subgroup, mk-def:identity_component_group_scheme, mk-lem:identity_component_geometrically_connected}
  Following Kleiman's proof of Lem.~5.1~(3) (subgroup part); the
  verbatim source quote is reproduced at
  \ref{mk-thm:identity_component_open_subgroup}.
  The product \(G^0 \times_k G^0\) is connected (EGA IV\(_2\) 4.5.8):
  \(G^0\) is geometrically connected
  (\ref{mk-lem:identity_component_geometrically_connected}) and every
  morphism to \(\Spec k\) is universally open, so the fibre product over
  \(k\) of the connected \(G^0\) with itself is connected. Its image under
  the multiplication \(\mu : G \times_k G \to G\) is therefore a connected
  subset of \(|G|\) containing \(\mu(e, e) = e\), hence contained in the
  clopen identity component \(G^0\)
  (\ref{mk-thm:identity_component_open_subgroup}). Likewise the image of the
  connected set \(G^0\) under the inversion of \(G\) is connected and
  contains \(e^{-1} = e\), hence lies in \(G^0\). Consequently, for every
  \(k\)-scheme \(T\), the subset of \(\operatorname{Hom}_k(T, G)\)
  consisting of morphisms whose set-theoretic image lies in \(G^0\) is a
  subgroup, functorially in \(T\); and since a morphism lands in the open
  subscheme \(G^0\) exactly when its set-theoretic image lies in \(G^0\),
  this subgroup presheaf is represented by \(G^0\). By the Yoneda
  characterisation of group objects, \(G^0\) inherits a \(k\)-group-scheme
  structure from \(G\), for which the inclusion \(G^0 \hookrightarrow G\)
  is a homomorphism of \(k\)-group schemes.
\end{proof}

\begin{theorem}[Identity component is of finite type and geometrically irreducible]\label{mk-thm:identity_component_finite_type_geom_irreducible}
  \dcref{5.1,custom}

  \textit{Source: \cite{kleiman-picard}, ``The Picard scheme'', \S 5,
  Lem.~5.1~(3); the verbatim source quote of (3) is at
  \ref{mk-thm:identity_component_open_subgroup}.}
  Let \(k\) be a field and \(G\) a \(k\)-group scheme locally of finite
  type. Then the identity component \(G^0\) is of finite type over
  \(k\) --- equivalently, locally of finite type \emph{and}
  quasi-compact --- and geometrically irreducible: after base change
  to the algebraic closure \(\bar k\), the underlying scheme
  \(G^0_{\bar k}\) is an irreducible \(\bar k\)-scheme.
\end{theorem}

\begin{proof}\uses{mk-thm:identity_component_open_subgroup, mk-thm:identity_component_base_change_commutes, mk-thm:identity_component_is_subgroup_homomorphism, mk-def:identity_component_group_scheme}
  Following Kleiman's proof of Lem.~5.1~(3) (finite-type and
  geometric-irreducibility part); the verbatim source quote is
  reproduced at \ref{mk-thm:identity_component_open_subgroup}. Base-change to
  \(\overline k\) is permitted because formation of \(G^0\) commutes with
  field extension (\ref{mk-thm:identity_component_base_change_commutes}), and
  over \(\overline k\) the identity component \(G^0\) carries the group
  structure inherited from \(G\)
  (\ref{mk-thm:identity_component_is_subgroup_homomorphism}).

  \emph{Irreducibility over \(\overline k\).}  Inside a Noetherian affine
  open of \(G^0\) (Noetherian because \(G\) is locally of finite type over a
  field), removing the union of all irreducible components but one leaves a
  nonempty irreducible open (Stacks 0052 (3)); its image under the
  inclusion of the affine open is a nonempty irreducible open
  \(U \subseteq G^0\). Pick a closed point \(g \in U\) (\(G^0\) is a
  Jacobson space, being locally of finite type over a field). For every
  closed point \(z\) of \(G^0\), \(g^{-1}z\) is again a closed --- hence,
  \(\overline k\) being algebraically closed, rational --- point, and right
  translation \(x \mapsto x \cdot (g^{-1}z)\) is an automorphism of \(G^0\)
  carrying \(g\) to \(z\); it therefore carries \(U\) to an irreducible open
  neighbourhood of \(z\). Consequently the union of all irreducible opens of
  \(G^0\) contains every closed point; since this union is open, its
  complement is closed, and a nonempty closed subset of a Jacobson space
  contains a closed point (Jacobson density) --- which cannot lie outside
  every constructed neighbourhood --- the complement is empty, so the union
  of irreducible opens is all of \(G^0\). By EGA I 6.1.10 (a connected space
  in which every point has an irreducible open neighbourhood is
  irreducible) and connectedness of \(G^0\)
  (\ref{mk-thm:identity_component_open_subgroup}), \(G^0\) is irreducible.

  \emph{Quasi-compactness over \(\overline k\).}  Fix an affine open
  \(W \ni e\). For a closed point \(z\), the ``division'' morphism
  \(\chi_z : x \mapsto x^{-1}z\) satisfies \(\chi_z(z) = e \in W\), so
  \(W \cap \chi_z^{-1}(W)\) is a nonempty open subset of the irreducible
  \(G^0\), hence contains a closed point \(g\) by Jacobson density. Then
  \(u := g^{-1}z\) is a rational point of \(W\) and \(z = g \cdot u\) lies in
  the image of the multiplication map \(m : W \times_{\overline k} W \to
  G^0\). Thus every closed point of \(G^0\) lies in the image of \(m\),
  which is quasi-compact as the image of the quasi-compact scheme
  \(W \times_{\overline k} W\); a finite affine subcover of this image,
  together with Jacobson density, covers all of \(G^0\), so \(G^0\) is
  quasi-compact.

  \emph{Descent to \(k\).}  Both conjuncts descend along the surjective
  projection \((G^0)_{\overline k} \to G^0\) coming from the base-change
  identification \((G_{\overline k})^0 \cong (G^0)_{\overline k}\)
  (\ref{mk-thm:identity_component_base_change_commutes}): a scheme admitting a
  surjection from a quasi-compact (resp. irreducible) scheme is itself
  quasi-compact (resp. irreducible). Quasi-compactness together with the
  standing locally-of-finite-type hypothesis on \(G\) makes \(G^0\) of
  finite type over \(k\). For geometric irreducibility at an arbitrary field
  extension \(K/k\), apply the same descent along the surjective projection
  from the algebraic closure of \(K\).
\end{proof}

\begin{theorem}[Formation of the identity component commutes with base change]\label{mk-thm:identity_component_base_change_commutes}
  \dcref{5.1,custom}

  \uses{mk-thm:identity_component_open_subgroup, mk-lem:geometricallyConnected_of_connected_of_section, mk-lem:identity_component_geometrically_connected, mk-def:identity_component_group_scheme}
  \textit{Source: \cite{kleiman-picard}, ``The Picard scheme'', \S 5,
  Lem.~5.1~(3); the verbatim source quote of (3) is at
  \ref{mk-thm:identity_component_open_subgroup}.}
  Let \(k\) be a field, \(G\) a \(k\)-group scheme locally of finite
  type, and \(K/k\) a field extension. The natural map
  \((G^0)_K \to (G_K)^0\), induced by base-changing the inclusion of
  \ref{mk-thm:identity_component_open_subgroup} along \(\Spec K \to \Spec k\)
  and using the universal property of the identity component on
  \(G_K\), is an isomorphism of \(K\)-schemes.
\end{theorem}

\begin{proof}
  \uses{mk-thm:identity_component_open_subgroup, mk-lem:geometricallyConnected_of_connected_of_section, mk-lem:identity_component_geometrically_connected, mk-def:identity_component_group_scheme}
  Following Kleiman's proof of Lem.~5.1~(3) (base-change part); the
  verbatim source quote is reproduced at
  \ref{mk-thm:identity_component_open_subgroup}.
  Geometric connectedness of \(G^0\) follows from the existence of the
  \(k\)-rational point \(e \in G^0(k)\): the identity component is
  connected and carries a \(k\)-rational section, so it is geometrically
  connected by
  \ref{mk-lem:geometricallyConnected_of_connected_of_section} (the
  Stacks 037Q criterion; cf.\ EGA IV\(_2\) 4.5.14). In particular for
  any field extension \(K/k\) the base change \((G^0)_K\) is connected,
  contains the image of the identity section of \(G_K\), and is open and
  closed in \(G_K\) (since open- and closed-immersion both descend under
  base change, \ref{mk-thm:identity_component_open_subgroup}). Hence
  \((G^0)_K\) coincides, as a subscheme of \(G_K\), with the identity
  component \((G_K)^0\) of \(G_K\); the natural comparison map
  \((G^0)_K \to (G_K)^0\) is therefore an isomorphism.
\end{proof}

\subsection{Descent of connectedness along field extensions}
\label{mk-subsec:connectedness_descent}

The bridge from connectedness of \(|G^0|\) to geometric connectedness of
\(G^0 \to \Spec k\) is the classical statement that a connected
\(k\)-scheme with a \(k\)-rational point is geometrically connected
(Stacks 04KV, EGA IV\(_2\) 4.5.14). We prove it from first principles.
The transcendental heart is the fact that over an algebraically closed
field the tensor product of two field extensions is a domain
(Stacks 05P3); the geometric assembly combines universal openness of
morphisms to the spectrum of a field (Stacks 0383, available in
mathlib), closedness of integral base changes, and a clopen-image
argument at the rational point.

\begin{lemma}[Evaluation at a maximal ideal over an algebraically closed field]\label{mk-lem:maximal_ideal_evaluation_alg_closed}
  \dcref{00FV,custom}

  \textit{Source: Zariski's lemma; Stacks, tag 00FV.}
  Let \(k\) be an algebraically closed field, \(R\) a finite-type
  \(k\)-algebra, and \(\mathfrak m \subseteq R\) a maximal ideal. Then
  there is a \(k\)-algebra homomorphism \(\varphi : R \to k\) with
  \(\ker \varphi = \mathfrak m\).
\end{lemma}
\begin{proof}
  The quotient \(R/\mathfrak m\) is a field and a finite-type
  \(k\)-algebra, hence a finite (in particular integral) extension of
  \(k\) by Zariski's lemma. Since \(k\) is algebraically closed, the
  structure map \(k \to R/\mathfrak m\) is bijective. Composing the
  quotient map \(R \to R/\mathfrak m\) with its inverse yields
  \(\varphi\), whose kernel is \(\mathfrak m\).
\end{proof}

\begin{lemma}[Tensor product with a finite-type domain over an algebraically closed field]\label{mk-lem:tensor_domain_finite_type_alg_closed}
  \dcref{05P3,custom}

  \uses{mk-lem:maximal_ideal_evaluation_alg_closed}
  \textit{Source: Stacks, tag 05P3 (finite-type case).}
  Let \(k\) be an algebraically closed field, \(R\) a finite-type
  \(k\)-algebra which is a domain, and \(A/k\) a field extension. Then
  \(R \otimes_k A\) is a domain.
\end{lemma}
\begin{proof}
  \uses{mk-lem:maximal_ideal_evaluation_alg_closed}
  Nontriviality holds because \(R\) and \(A\) are nontrivial and embed
  their fraction fields compatibly. Fix a \(k\)-basis \((a_i)\) of
  \(A\); then \(R \otimes_k A\) is a free \(R\)-module with basis
  \((1 \otimes a_i)\), so every element \(x\) has unique coordinates
  \(c_i(x) \in R\). Suppose \(x y = 0\). For every maximal ideal
  \(\mathfrak m \subseteq R\), choose
  \(\varphi : R \to k\) with kernel \(\mathfrak m\) by
  \ref{mk-lem:maximal_ideal_evaluation_alg_closed}, and apply the ring
  homomorphism \(\varphi \otimes \id : R \otimes_k A \to A\): it sends
  \(x \mapsto \sum_i \varphi(c_i(x)) a_i\). Since \(A\) is a field,
  \((\varphi \otimes \id)(x) = 0\) or \((\varphi \otimes \id)(y) = 0\);
  by linear independence of \((a_i)\) over \(k\) this means all
  \(c_i(x) \in \mathfrak m\) or all \(c_j(y) \in \mathfrak m\). In
  either case \(c_i(x)\, c_j(y) \in \mathfrak m\) for all \(i, j\) and
  all maximal \(\mathfrak m\). A finite-type algebra over a field is a
  Jacobson ring, and in a Jacobson domain the intersection of all
  maximal ideals is zero; hence \(c_i(x)\, c_j(y) = 0\) for all
  \(i, j\). If \(x \neq 0\), some \(c_{i_0}(x) \neq 0\), and since
  \(R\) is a domain every \(c_j(y) = 0\), i.e.\ \(y = 0\).
\end{proof}

\begin{lemma}[Tensor product of field extensions over an algebraically closed field]\label{mk-lem:tensor_domain_alg_closed}
  \dcref{05P3,custom}

  \uses{mk-lem:tensor_domain_finite_type_alg_closed}
  \textit{Source: Stacks, tag 05P3.}
  Let \(k\) be an algebraically closed field and \(A/k\), \(B/k\) field
  extensions. Then \(A \otimes_k B\) is a domain.
\end{lemma}
\begin{proof}
  \uses{mk-lem:tensor_domain_finite_type_alg_closed}
  Suppose \(x y = 0\) with \(x, y \in A \otimes_k B\). Writing \(x\) and
  \(y\) as finite sums of elementary tensors, let \(R \subseteq A\) be
  the \(k\)-subalgebra generated by the finitely many first components;
  \(R\) is a finite-type \(k\)-algebra and a domain (a subring of the
  field \(A\)). Since \(B\) is free, hence flat, as a \(k\)-module, the
  natural map \(R \otimes_k B \to A \otimes_k B\) is injective. The
  elements \(x, y\) lift to \(x', y' \in R \otimes_k B\) with
  \(x' y' = 0\); by \ref{mk-lem:tensor_domain_finite_type_alg_closed}
  (applied with the field extension \(B\)) we get \(x' = 0\) or
  \(y' = 0\), hence \(x = 0\) or \(y = 0\). Nontriviality is as in
  \ref{mk-lem:tensor_domain_finite_type_alg_closed}.
\end{proof}

\begin{lemma}[The spectrum of a domain is connected]\label{mk-lem:spec_domain_connected}
  \dcref{custom}

  Let \(R\) be a commutative domain. Then \(\Spec R\) is a connected
  (indeed irreducible) nonempty topological space.
\end{lemma}
\begin{proof}
  In a domain the nilradical \((0)\) is prime, so \(\Spec R\) is
  irreducible with generic point \((0)\); irreducible spaces are
  connected and nonempty.
\end{proof}

\begin{lemma}[Projections from constant-field base changes are geometrically connected]\label{mk-lem:geometrically_connected_projection_alg_closed}
  \dcref{custom}

  \uses{mk-lem:tensor_domain_alg_closed, mk-lem:spec_domain_connected}
  Let \(k\) be an algebraically closed field, \(L/k\) a field extension,
  and \(f : X \to \Spec k\) a morphism of schemes. Then the projection
  \(q : X \times_k \Spec L \to X\) is geometrically connected: for every
  field \(K\) and morphism \(\Spec K \to X\), the fiber product
  \((X \times_k \Spec L) \times_X \Spec K\) is connected.
\end{lemma}
\begin{proof}
  \uses{mk-lem:tensor_domain_alg_closed, mk-lem:spec_domain_connected}
  By the pasting law for fiber products,
  \((X \times_k \Spec L) \times_X \Spec K \cong \Spec K \times_{\Spec k}
  \Spec L = \Spec (K \otimes_k L)\), where \(K\) is a field extension of
  \(k\) via \(\Spec K \to X \to \Spec k\). By
  \ref{mk-lem:tensor_domain_alg_closed} the ring \(K \otimes_k L\) is a
  domain, so its spectrum is connected by
  \ref{mk-lem:spec_domain_connected}.
\end{proof}

\begin{lemma}[Connectedness is stable under base change from an algebraically closed field]\label{mk-lem:connected_base_change_alg_closed}
  \dcref{0363,custom}

  \uses{mk-lem:geometrically_connected_projection_alg_closed}
  \textit{Source: Stacks, tag 0363.}
  Let \(k\) be an algebraically closed field, \(L/k\) a field extension,
  and \(f : X \to \Spec k\) a morphism with \(|X|\) connected. Then
  \(X \times_k \Spec L\) is connected.
\end{lemma}
\begin{proof}
  \uses{mk-lem:geometrically_connected_projection_alg_closed}
  The projection \(q : X \times_k \Spec L \to X\) is open: it is the
  base change of \(\Spec L \to \Spec k\), which is universally open
  because every morphism to the spectrum of a field is universally open
  (Stacks 0383; available in mathlib). By
  \ref{mk-lem:geometrically_connected_projection_alg_closed} the
  fibers of \(q\) over points of \(X\) are connected. An open
  surjective map with connected fibers onto a connected space has
  connected total space (mathlib,
  \(\texttt{GeometricallyConnected.connectedSpace}\)).
\end{proof}

\begin{lemma}[Base change to the algebraic closure preserves connectedness given a section]\label{mk-lem:connected_base_change_algebraic_closure}
  \dcref{04KV,custom}

  \textit{Source: Stacks, tag 04KV (algebraic step); EGA IV\(_2\) 4.5.13.}
  Let \(k\) be a field, \(f : X \to \Spec k\) a morphism with \(|X|\)
  connected admitting a section \(s : \Spec k \to X\), and
  \(\bar k\) an algebraic closure of \(k\). Then
  \(X \times_k \Spec \bar k\) is connected.
\end{lemma}
\begin{proof}
  The projection \(p : X_{\bar k} := X \times_k \Spec \bar k \to X\) is:
  open (base change of the universally open \(\Spec \bar k \to \Spec
  k\), as in \ref{mk-lem:connected_base_change_alg_closed}); closed
  (base change of \(\Spec \bar k \to \Spec k\), which is integral
  because \(\bar k / k\) is algebraic, and integral morphisms are
  universally closed); and surjective (base change of a surjection).
  Let \(x_0 = s(\ast)\) be the image of the section. The section
  induces a \(k\)-algebra retraction \(\kappa(x_0) \to k\) of the field
  map \(k \to \kappa(x_0)\); a field homomorphism admitting a
  retraction is bijective, so \(\kappa(x_0) = k\). Hence the fiber
  \(p^{-1}(x_0) = \Spec(\kappa(x_0) \otimes_k \bar k) = \Spec \bar k\)
  is a single point \(\zeta\).

  \(X_{\bar k}\) is nonempty (it contains \(\zeta\), or equivalently
  the image of the base-changed section). Let \(T \subseteq
  X_{\bar k}\) be clopen with \(T \neq \emptyset\) and
  \(T^{\mathsf c} \neq \emptyset\). Since \(p\) is both open and
  closed, \(p(T)\) and \(p(T^{\mathsf c})\) are nonempty clopen subsets
  of the connected space \(|X|\), hence both equal \(|X|\). In
  particular \(x_0 \in p(T) \cap p(T^{\mathsf c})\), so the fiber
  \(p^{-1}(x_0)\) meets both \(T\) and \(T^{\mathsf c}\); but that
  fiber is the single point \(\zeta\) — a contradiction. So
  \(X_{\bar k}\) admits no nontrivial clopen partition and is
  connected.
\end{proof}

\begin{theorem}[Connected with a rational section implies geometrically connected]\label{mk-thm:geometrically_connected_of_section}
  \dcref{04KV,custom}

  \uses{mk-lem:connected_base_change_algebraic_closure, mk-lem:connected_base_change_alg_closed}
  \textit{Source: Stacks, tag 04KV; EGA IV\(_2\) 4.5.14.}
  Let \(k\) be a field and \(f : X \to \Spec k\) a morphism with
  \(|X|\) connected, admitting a section \(s : \Spec k \to X\). Then
  \(f\) is geometrically connected: for every field extension \(L/k\)
  the base change \(X \times_k \Spec L\) is connected.
\end{theorem}
\begin{proof}
  \uses{mk-lem:connected_base_change_algebraic_closure, mk-lem:connected_base_change_alg_closed}
  Let \(L/k\) be arbitrary, \(\bar L\) an algebraic closure of \(L\),
  and \(\bar k\) an algebraic closure of \(k\). Since \(\bar k/k\) is
  algebraic and \(\bar L\) is algebraically closed, there is a
  \(k\)-embedding \(\bar k \hookrightarrow \bar L\). By
  \ref{mk-lem:connected_base_change_algebraic_closure},
  \(X_{\bar k}\) is connected; by
  \ref{mk-lem:connected_base_change_alg_closed} applied over the
  algebraically closed field \(\bar k\) to the extension
  \(\bar L / \bar k\), the scheme
  \(X_{\bar L} = X_{\bar k} \times_{\bar k} \Spec \bar L\) is
  connected. Finally the projection \(X_{\bar L} \to X_L\) is
  surjective (base change of the surjection
  \(\Spec \bar L \to \Spec L\)) and continuous, so \(X_L\) is the
  continuous image of a connected space, hence connected.
\end{proof}
\begin{lemma}[Stacks 037Q iff-direction: connected + section ⟹ geometrically connected]\label{mk-lem:geometricallyConnected_of_connected_of_section}
  \dcref{037Q,custom}

  \uses{mk-thm:geometrically_connected_of_section}
  \textit{Source: Stacks, tag 037Q (iff-direction (2)\(\Rightarrow\)(1)).}
  Let \(k\) be a field and let \(f : X \to \Spec k\) be a morphism of
  schemes such that the underlying topological space \(|X|\) is
  connected. Suppose there is a \(k\)-rational section \(s : \Spec k \to X\)
  of \(f\) (so \(s \circ f = \id_X\) topologically and the image of \(s\)
  meets every irreducible component of \(X\) through which \(s\) factors).
  Then \(f : X \to \Spec k\) is geometrically connected, in the sense
  that for every field extension \(K/k\) the base change \(X_K\) is
  connected.
\end{lemma}
\begin{proof}
  \uses{mk-thm:geometrically_connected_of_section}
  Immediate from \ref{mk-thm:geometrically_connected_of_section} applied
  to \(f\) and the section \(s\).
\end{proof}

\subsection{Topological and connectedness substrate (helper lemmas)}
\label{mk-subsec:identity_component_substrate}


\begin{lemma}
[Finitely many connected components on a Noetherian space]
  \label{mk-lem:noetherian_finite_connected_components}
  \dcref{custom}

  Let \(\alpha\) be a Noetherian topological space. Then the set of
  connected components of \(\alpha\) is finite.
\end{lemma}
\begin{proof}
  The assignment sending an irreducible component \(C\) of \(\alpha\) to the
  connected component of an arbitrarily chosen point of \(C\) is
  well-defined (an irreducible component is preconnected, hence lies in a
  single connected component) and surjective (every point lies in the
  irreducible component through it, which is contained in its connected
  component). A Noetherian space has finitely many irreducible components,
  so the connected components are finite as the image of a finite set.

\end{proof}

\begin{lemma}[Connected components are open on a Noetherian space]\label{mk-lem:noetherian_isopen_connected_component}
  \dcref{custom}

  \uses{mk-lem:noetherian_finite_connected_components}
  Let \(\alpha\) be a Noetherian topological space and \(x \in \alpha\). Then
  the connected component \(\connComp(x)\) of \(x\) is open in \(\alpha\).
\end{lemma}
\begin{proof}
  \uses{mk-lem:noetherian_finite_connected_components}
  The space of connected components of \(\alpha\) is totally disconnected,
  hence \(T_1\), and finite by
  \ref{mk-lem:noetherian_finite_connected_components}; a finite \(T_1\) space
  carries the discrete topology, so each singleton is open. The preimage of
  the singleton \(\{\,[x]\,\}\) under the continuous quotient map
  \(\alpha \to \pi_0(\alpha)\) is exactly \(\connComp(x)\), which is
  therefore open.
\end{proof}

\begin{lemma}[The underlying space of a locally-finite-type \(k\)-scheme is locally connected]\label{mk-lem:identity_component_locally_connected}
  \dcref{custom}

  \uses{mk-lem:noetherian_isopen_connected_component}
  Let \(k\) be a field and \(G\) an object of the over-category over
  \(\Spec k\) whose structural morphism \(G \to \Spec k\) is locally of
  finite type. Then the underlying topological space \(|G|\) is locally
  connected.
\end{lemma}
\begin{proof}
  \uses{mk-lem:noetherian_isopen_connected_component}
  Being locally of finite type over the field \(k\), the scheme \(G\) is
  locally Noetherian. Each point \(y \in |G|\) has an affine open
  neighbourhood \(W = \Spec R\) with \(R\) Noetherian, so \(|W|\) is a
  Noetherian space; by
  \ref{mk-lem:noetherian_isopen_connected_component} the connected component
  of \(y\) inside \(W\) is open, and its image under the open inclusion
  \(W \hookrightarrow |G|\) is an open connected neighbourhood of \(y\)
  contained in any prescribed open set through \(y\). This is precisely the
  local-connectedness criterion.
\end{proof}

\begin{definition}[The identity point]\label{mk-def:identity_section_point}
  \dcref{custom}

  Let \(k\) be a field and \(G\) a \(k\)-group scheme. The
  \emph{identity point} of \(G\) is the image \(e(\ast) \in |G|\) of the
  unique point of \(\Spec k\) under the identity section
  \(e : \Spec k \to G\) (well-defined because \(|\Spec k|\) is a
  one-point space).
\end{definition}
\begin{proof}
  section).
\end{proof}

\begin{definition}[The carrier of the identity component]\label{mk-def:identity_component_carrier}
  \dcref{custom}

  \uses{mk-def:identity_section_point, mk-lem:identity_component_locally_connected}
  Let \(k\) be a field and \(G\) a \(k\)-group scheme locally of finite
  type. The \emph{carrier of the identity component} is the open subset of
  \(|G|\) given by the connected component
  \(\connComp(e(\ast))\) of the identity point
  (\ref{mk-def:identity_section_point}), packaged as an open of \(G\); its
  openness is supplied by the local connectedness of \(|G|\)
  (\ref{mk-lem:identity_component_locally_connected}).
\end{definition}
\begin{proof}
  \uses{mk-def:identity_section_point, mk-lem:identity_component_locally_connected}
  The carrier set is \(\connComp(e(\ast))\); a connected component of a
  locally connected space is open
  (\ref{mk-lem:identity_component_locally_connected}). Proved directly in
  the formal proof.
\end{proof}

\begin{lemma}[The carrier of the identity component is connected]\label{mk-lem:identity_component_carrier_connected}
  \dcref{custom}

  \uses{mk-def:identity_component_carrier, mk-def:identity_section_point}
  Let \(k\) be a field and \(G\) a \(k\)-group scheme locally of finite
  type. The carrier of the identity component
  (\ref{mk-def:identity_component_carrier}), with its subspace topology, is a
  connected topological space.
\end{lemma}
\begin{proof}
  \uses{mk-def:identity_component_carrier, mk-def:identity_section_point}
  As a subspace, the carrier \(\connComp(e(\ast))\) is preconnected (a
  connected component is preconnected) and nonempty (it contains the
  identity point \(e(\ast)\), \ref{mk-def:identity_section_point}), hence
  connected.
\end{proof}

\begin{lemma}[The identity component is connected]\label{mk-lem:identity_component_connected}
  \dcref{custom}

  \uses{mk-lem:identity_component_carrier_connected, mk-def:identity_component_group_scheme}
  Let \(k\) be a field and \(G\) a \(k\)-group scheme locally of finite
  type. The underlying topological space of the identity component
  \(G^0\) (\ref{mk-def:identity_component_group_scheme}) is connected.
\end{lemma}
\begin{proof}
  \uses{mk-lem:identity_component_carrier_connected, mk-def:identity_component_group_scheme}
  By construction the underlying space of \(G^0\) coincides with the
  carrier of \ref{mk-def:identity_component_carrier}, whose connectedness is
  \ref{mk-lem:identity_component_carrier_connected}. Proved directly in
  the formal proof.
\end{proof}

\begin{lemma}[The identity section lands in the identity-component carrier]\label{mk-lem:identity_component_section_range_subset}
  \dcref{custom}

  \uses{mk-def:identity_component_carrier, mk-def:identity_section_point}
  Let \(k\) be a field and \(G\) a \(k\)-group scheme locally of finite
  type. The image of the identity section \(e : \Spec k \to G\) is
  contained in the carrier of the identity component
  (\ref{mk-def:identity_component_carrier}).
\end{lemma}
\begin{proof}
  \uses{mk-def:identity_component_carrier, mk-def:identity_section_point}
  The image of \(e\) is the single point \(e(\ast)\)
  (\ref{mk-def:identity_section_point}), which lies in its own connected
  component, i.e.\ in the carrier.
\end{proof}

\begin{definition}[The identity section into the identity component]\label{mk-def:identity_component_section}
  \dcref{custom}

  \uses{mk-def:identity_component_carrier, mk-lem:identity_component_section_range_subset, mk-def:identity_component_group_scheme}
  Let \(k\) be a field and \(G\) a \(k\)-group scheme locally of finite
  type. The identity section \(e : \Spec k \to G\) factors through the
  open immersion \(G^0 \hookrightarrow G\)
  (\ref{mk-def:identity_component_carrier}); the resulting morphism
  \(\Spec k \to G^0\) is the \emph{identity section into the identity
  component}, obtained by the universal property of the open immersion
  using the range containment of
  \ref{mk-lem:identity_component_section_range_subset}.
\end{definition}
\begin{proof}
  \uses{mk-def:identity_component_carrier, mk-lem:identity_component_section_range_subset}
  The lift exists and is unique by the universal property of the open
  immersion \(G^0 \hookrightarrow G\), the hypothesis being the range
  containment of \ref{mk-lem:identity_component_section_range_subset}.

\end{proof}

\begin{lemma}[The lifted identity section is a section]\label{mk-lem:identity_component_section_is_section}
  \dcref{custom}

  \uses{mk-def:identity_component_section, mk-lem:identity_component_section_range_subset, mk-def:identity_component_group_scheme}
  Let \(k\) be a field and \(G\) a \(k\)-group scheme locally of finite
  type. The morphism \(\Spec k \to G^0\) of
  \ref{mk-def:identity_component_section} is a section of the structural
  morphism \(G^0 \to \Spec k\): composing it with \(G^0 \to \Spec k\)
  recovers the identity on \(\Spec k\).
\end{lemma}
\begin{proof}
  \uses{mk-def:identity_component_section, mk-lem:identity_component_section_range_subset}
  Composing through the open immersion recovers the original identity
  section \(e\) (the defining factorisation of
  \ref{mk-def:identity_component_section}), and \(e\) is a section of
  \(G \to \Spec k\); hence the lift is a section of \(G^0 \to \Spec k\).

\end{proof}

\begin{lemma}[The identity component is geometrically connected]\label{mk-lem:identity_component_geometrically_connected}
  \dcref{custom}

  \uses{mk-lem:geometricallyConnected_of_connected_of_section, mk-def:identity_component_section, mk-lem:identity_component_section_is_section, mk-lem:identity_component_connected}
  Let \(k\) be a field and \(G\) a \(k\)-group scheme locally of finite
  type. The structural morphism \(G^0 \to \Spec k\) of the identity
  component is geometrically connected.
\end{lemma}
\begin{proof}
  \uses{mk-lem:geometricallyConnected_of_connected_of_section, mk-def:identity_component_section, mk-lem:identity_component_section_is_section, mk-lem:identity_component_connected}
  The identity component \(G^0\) is connected
  (\ref{mk-lem:identity_component_connected}) and carries the \(k\)-rational
  section of \ref{mk-def:identity_component_section}, which is a section of
  \(G^0 \to \Spec k\) by \ref{mk-lem:identity_component_section_is_section}.
  The criterion ``connected \(+\) \(k\)-rational section \(\Rightarrow\)
  geometrically connected''
  (\ref{mk-lem:geometricallyConnected_of_connected_of_section}) then yields
  the conclusion.
\end{proof}

\section{The identity component of the Picard scheme}
\label{mk-sec:pic0_definition}


\begin{definition}[The identity component of the Picard scheme]\label{mk-def:pic_zero_subscheme}
  \dcref{5.3,custom}

  \uses{mk-def:identity_component_group_scheme, mk-def:pic_scheme,
    mk-def:pic_scheme_lft, mk-thm:pic_is_group_scheme}
  \textit{Source: \cite{kleiman-picard}, ``The Picard scheme'', \S 5, opening +
  Prp.~5.3.}
  Let \(C/k\) be a smooth proper geometrically integral curve over a field
  \(k\), and assume the legacy relative Picard functor
  \(\Pic^\sharp_{C/k}\) is represented by \(\Pic_{C/k}\). The
  \emph{identity component of the Picard scheme}
  \(\Pic^0_{C/k}\) is the identity component
  \((\Pic_{C/k})^0\) of the \(k\)-group scheme \(\Pic_{C/k}\) of
  \ref{mk-def:pic_scheme}, in the sense of
  \ref{mk-def:identity_component_group_scheme}. By
  \ref{mk-thm:identity_component_open_subgroup}, \(\Pic^0_{C/k}\) is an open
  and closed subgroup scheme of \(\Pic_{C/k}\) of finite type over \(k\),
  geometrically irreducible, and its formation commutes with
  extension of the base field.
\end{definition}

\section{The degree map}
\label{mk-sec:pic_degree}

Return now to the arbitrary-field representer \(\Pic^{\et}_{C/k}\) of
\ref{mk-thm:fga_pic_representability_etale}. Its disjoint open pieces are indexed
by the Hilbert polynomial \(\Phi\) used in the FGA construction. Given a
\(T\)-point, choose an \'etale cover \(U\to T\) on which it is represented by
an invertible sheaf \(\mathcal L\) on \(C\times_kU\). With \(\mathcal H\) the
chosen very ample sheaf, the construction uses the fibrewise dual convention
\[
  \Phi_{\mathcal L,u}(n)
  =\chi(C_u,\mathcal L_u^{-1}\otimes\mathcal H_u^{\otimes n})
  =-\deg\mathcal L_u+n\deg\mathcal H_u+1-g
  \qquad (u\in U).
\]
Thus its coefficient of \(n\) is \(\deg\mathcal H_u\), while
\(\deg\mathcal L_u=-(\Phi_{\mathcal L,u}(0)-
\chi(C_u,\mathcal O_{C_u}))\). Intrinsically,
\(\deg\mathcal L_u=\chi(C_u,\mathcal L_u)-
\chi(C_u,\mathcal O_{C_u})\), equivalently the coefficient of \(n\) in
\(n\mapsto\chi(C_u,\mathcal L_u^{\otimes n})\).

\begin{definition}[Degree of a Picard family]\label{mk-def:divisor_degree_pic}
  \dcref{custom}
  \uses{mk-def:pic_et_sheaf, mk-def:hilbert_polynomial,
    mk-thm:fga_pic_representability_etale}
  \textit{Source: \cite{milne-av}, ``Abelian Varieties'' (course notes, 2008),
  \S III.1, p.~88.}
  Let \(C/k\) be a smooth proper geometrically integral curve, and let
  \(\Pic^{\et}_{C/k}\) represent the \'etale-sheafified relative Picard
  functor. For every \(k\)-scheme \(T\), the \emph{family degree} is an
  additive map
  \[
    \deg_T:\Pic_{(C/k)\et}(T)
      \longrightarrow \operatorname{LC}(|T|,\mathbb Z),
  \]
  natural under pullback in \(T\). Here
  \(\operatorname{LC}(|T|,\mathbb Z)\) denotes the additive group of locally
  constant integer-valued functions on the underlying space. If a class is
  represented after an \'etale cover \(U\to T\) by \(\mathcal L\) on
  \(C\times_kU\), its pullback has value at \(u\in U\)
  \[
    \deg_U(\mathcal L)(u)=\deg(\mathcal L_u)
    =\chi(C_u,\mathcal L_u)-\chi(C_u,\mathcal O_{C_u}).
  \]
  These fibre values are locally constant on \(U\), invariant under the
  relative Picard quotient and \'etale descent relation, and therefore define
  \(\deg_T\). For \(f:T'\to T\), one has
  \(\deg_{T'}(f^*\lambda)=\deg_T(\lambda)\circ|f|\); fibrewise additivity of
  divisor degree gives additivity.

  Taking \(T=\Spec k\) yields
  \(\deg:\Pic^{\et}_{C/k}(k)\to\mathbb Z\). The subfunctor on which
  \(\deg_T\) is constantly \(d\) is represented by an open-and-closed degree
  piece \((\Pic^{\et}_{C/k})^d\), giving
  \(\Pic^{\et}_{C/k}=\coprod_{d\in\mathbb Z}(\Pic^{\et}_{C/k})^d\). After
  base change to \(\overline k\), these are the geometric connected components.
\end{definition}

\section{The Pic-zero abelian variety}
\label{mk-sec:pic0_abelian_variety}

The terminal statement of the chapter identifies \(\Pic^0_{C/k}\) with
an abelian variety of dimension \(g\) --- the Jacobian variety of \(C\).
Recall (Milne~\S I.1, p.~8) that a \(k\)-group scheme is an
\emph{abelian variety} iff it is a complete (proper) connected group
variety over \(k\); commutativity and projectivity then follow
automatically.

\begin{theorem}[\(\Pic^0_{C/k}\) is an abelian variety]\label{mk-thm:pic_zero_is_abelian_variety}
  \dcref{1.4,custom}

  \uses{mk-def:pic_zero_subscheme, mk-thm:identity_component_open_subgroup, mk-thm:identity_component_is_subgroup_homomorphism, mk-thm:identity_component_finite_type_geom_irreducible, mk-thm:identity_component_base_change_commutes, mk-thm:fga_pic_representability, mk-thm:pic_is_group_scheme}
  \textit{Source: \cite{kleiman-picard}, ``The Picard scheme'', \S 5,
  Ex.~5.23 + Rmk.~5.26; cf.\ \cite{milne-av}, ``Abelian Varieties'',
  \S I.1 (def.\ of abelian variety, p.~8).}
  Let \(C/k\) be a smooth proper geometrically integral curve of genus
  \(g\) over a field \(k\). The identity component \(\Pic^0_{C/k}\) of
  \ref{mk-def:pic_zero_subscheme} is an \emph{abelian variety over
  \(k\)}: it is proper, smooth, and geometrically irreducible as a
  \(k\)-scheme, and it carries a \(k\)-group-scheme structure. (In the
  \texttt{[IsProper]}, \texttt{[Smooth]}, and
  \texttt{[GeometricallyIrreducible]} on the structural morphism
  \texttt{(Pic0Scheme C).hom} together with the existence of a
  \texttt{[GrpObj]} structure on \(\Pic^0_{C/k}\). Commutativity
  follows automatically from Milne~\S I.1, Cor.~1.4 and is not
  separately stated.)
\end{theorem}
\begin{proof}
  \uses{mk-def:pic_zero_subscheme, mk-thm:identity_component_open_subgroup, mk-thm:identity_component_is_subgroup_homomorphism, mk-thm:identity_component_finite_type_geom_irreducible, mk-thm:identity_component_base_change_commutes, mk-thm:fga_pic_representability, mk-thm:pic_is_group_scheme, mk-thm:pic0_isAbelianVariety}
  We assemble the four conjuncts of the abelian-variety structure on
  \(\Pic^0_{C/k}\) (\ref{mk-def:pic_zero_subscheme}) from four inputs.

  (i) \emph{Clopen subgroup of finite type, geometrically irreducible.}
  By
  \ref{mk-thm:identity_component_open_subgroup},
  \ref{mk-thm:identity_component_is_subgroup_homomorphism},
  \ref{mk-thm:identity_component_finite_type_geom_irreducible}, and
  \ref{mk-thm:identity_component_base_change_commutes}
  applied to \(G = \Pic_{C/k}\) (a \(k\)-group scheme locally of finite
  type by
  \ref{mk-thm:pic_is_group_scheme} and \ref{mk-thm:fga_pic_representability}),
  \(\Pic^0_{C/k}\) is an open and closed subgroup scheme of \(\Pic_{C/k}\),
  of finite type over \(k\), geometrically irreducible, and its
  formation commutes with extending \(k\). In particular \(\Pic^0_{C/k}\)
  is separated and locally of finite type as a subscheme of
  \(\Pic_{C/k}\), and connected.

  (ii) \emph{Quasi-projectivity.} The disjoint-union structure of
  \(\Pic_{C/k}\) (\ref{mk-thm:fga_pic_representability}: a disjoint union
  of open quasi-projective \(k\)-subschemes) restricted to \(\Pic^0_{C/k}\)
  expresses the identity component as a quasi-projective \(k\)-scheme
  (Kleiman~\S 5, Thm.~5.4: \(X/k\) projective and geometrically
  integral \(\Longrightarrow\) \(\Pic^0_{X/k}\) is quasi-projective).

  (iii) \emph{Projectivity (properness).} For a smooth proper
  geometrically integral curve \(C/k\) of genus \(g > 0\), \(C/k\) is
  also geometrically normal (the smoothness implies geometric
  regularity, hence geometric normality). Kleiman~\S 5, Thm.~5.4
  then upgrades the quasi-projective conclusion to projective:
  \(\Pic^0_{C/k}\) is projective over \(k\), in particular proper.
  The proof of Theorem~5.4 uses the Chevalley--Rosenlicht structure
  theorem to reduce to ruling out non-constant maps
  \(\mathbb G_m \to \Pic^0_{C/k}\), which on a geometrically normal
  \(X\) becomes a divisor-theoretic computation.

  (iv) \emph{Smoothness.} For a smooth proper curve \(C/k\), the
  Picard scheme is smooth at \(0\): this is Exercise~5.23 of
  Kleiman~\S 5, which asserts that for \(X/S\) projective and flat
  with integral geometric fibres of arithmetic genus \(p_a\), the
  generalised Jacobians \(\Pic^0_{X_s/k_s}\) form a smooth
  quasi-projective family of relative dimension \(p_a\); specialising
  to \(S = \Spec k\) and \(X = C\) a smooth proper curve of (geometric
  = arithmetic) genus \(g\) gives smoothness of \(\Pic^0_{C/k}\).

  Combining (i)--(iv): \(\Pic^0_{C/k}\) is a smooth, projective,
  geometrically irreducible \(k\)-group scheme. A proper connected
  \(k\)-group variety is by definition an \emph{abelian variety}
  (Milne~\S I.1, p.~8), and is automatically commutative (Milne~\S I.1,
  Cor.~1.4). This establishes
  \ref{mk-thm:pic_zero_is_abelian_variety}.
\end{proof}

\begin{theorem}[Dimension of \(\Pic^0_{C/k}\) equals the genus of \(C\)]\label{mk-thm:pic_zero_dimension_equals_genus}
  \dcref{5.13,custom}

  \uses{mk-thm:pic_zero_is_abelian_variety, mk-def:genus}
  \textit{Source: \cite{milne-av}, ``Abelian Varieties'' (course notes, 2008),
  \S III.1, Rmk.~1.4(e), p.~86; cf.\ \cite{kleiman-picard}, ``The Picard scheme'',
  \S 5, Cor.~5.13 + Ex.~5.23.}
  Let \(C/k\) be a smooth proper geometrically integral curve of genus
  \(g = g(C)\) (\ref{mk-def:genus}) over a field \(k\). The dimension of
  the abelian variety \(\Pic^0_{C/k}\)
  (\ref{mk-thm:pic_zero_is_abelian_variety}) equals the genus:
  \[
    \dim_k \Pic^0_{C/k} \;=\; g(C).
  \]
\end{theorem}
\begin{proof}
  \uses{mk-thm:pic_zero_is_abelian_variety, mk-def:genus}
  By Kleiman~\S 5,
  Cor.~5.13, the inequality
  \(\dim \Pic_{C/k} \le \dim_k H^1(C, \mathcal{O}_C)\) holds, with
  equality if and only if \(\Pic_{C/k}\) is smooth at \(0\), in which
  case \(\Pic_{C/k}\) is smooth of dimension \(\dim_k H^1(C,
  \mathcal{O}_C)\) everywhere. The abelian-variety structure of
  \ref{mk-thm:pic_zero_is_abelian_variety} (step (iv) of its proof)
  shows \(\Pic^0_{C/k}\) is smooth at \(0\); equality therefore holds
  on the open identity component, and
  \(\dim_k \Pic^0_{C/k} = \dim_k H^1(C, \mathcal{O}_C) = g(C)\)
  by \ref{mk-def:genus}. This is Milne~\S III.1, Rmk.~1.4(e): ``The
  dimension of \(J\) is the genus of \(C\)''.
\end{proof}

\begin{theorem}[The identity component of \(\Pic^{\et}_{C/k}\) is the kernel of degree on \(k\)-points]\label{mk-thm:pic_zero_k_points_iff_degree_zero}
  \dcref{custom}
  \uses{mk-thm:identity_component_open_subgroup, mk-def:divisor_degree_pic,
    mk-def:pic_et_sheaf, mk-thm:fga_pic_representability_etale}
  \textit{Source: \cite{milne-av}, ``Abelian Varieties'' (course notes, 2008),
  \S III.1, p.~88 (``\(\Pic^0(C)\)\dots\ degree zero''); the verbatim
  source quote is at \ref{mk-def:divisor_degree_pic}.}
  Let \(C/k\) be a smooth proper geometrically integral curve of genus
  \(g\) over a field \(k\), and let \(\Pic^{\et}_{C/k}\) represent
  \(\Pic_{(C/k)\et}\) as in \ref{mk-thm:fga_pic_representability_etale}. A
  \(k\)-point \(\lambda \in \Pic^{\et}_{C/k}(k)\) lies in the image of the
  identity-component inclusion
  \((\Pic^{\et}_{C/k})^0\hookrightarrow\Pic^{\et}_{C/k}\)
  of \ref{mk-thm:identity_component_open_subgroup}
  if and only if the degree
  \(\deg(\lambda) \in \mathbb{Z}\) of \ref{mk-def:divisor_degree_pic}
  vanishes. Equivalently,
  \((\Pic^{\et}_{C/k})^0(k)
    =\ker\bigl(\deg:\Pic^{\et}_{C/k}(k)\to\mathbb Z\bigr)\).
\end{theorem}

\begin{proof}
  \uses{mk-thm:identity_component_open_subgroup, mk-def:divisor_degree_pic}
  A class
  \(\lambda \in \Pic^{\et}_{C/k}(k)\) lies in the connected component of
  the identity --- i.e.\ in the image of the open immersion
  \((\Pic^{\et}_{C/k})^0\hookrightarrow\Pic^{\et}_{C/k}\) --- precisely
  when its Hilbert-polynomial stratum is the degree-zero stratum.
  On an \'etale cover carrying a representative \(\mathcal L\), this says
  fibrewise that
  \(\chi(C_u,\mathcal L_u)-\chi(C_u,\mathcal O_{C_u})=0\).
  Hence
  \((\Pic^{\et}_{C/k})^0(k)
    =\ker(\deg:\Pic^{\et}_{C/k}(k)\to\mathbb Z)\),
  which is the content of Milne~\S III.1, p.~88 (the verbatim quote
  at \ref{mk-def:divisor_degree_pic} explicitly writes ``\(\Pic^0(C)\)
  for the group of isomorphism classes of invertible sheaves of degree
  zero on \(C\)'').
\end{proof}
